The measurement of transverse energy per unit pseudorapidity (\detdeta) using data from Run 2024 in \auau collisions at \snn = 200 GeV performed with the sPHENIX detector over the range $\left|\eta\right| < 1.1$. 
The \detdeta is measured with the sPHENIX calorimeter system, including the use of a hadronic calorimeter to measure \detdeta for the first time at RHIC. 
The measurement is first performed using only the EMCal and only the HCal, with good agreement between the two despite the different sensitivities and calibration procedures for these detector systems. 
The measurement is performed with the full calorimeter system, differentially in centrality over the range $0$--$70$\%, and is also reported as a ratio to the estimated number of participant pairs (\Npart/2). 
The sPHENIX full calorimeter results are compared to previous STAR and PHENIX results, with sPHENIX providing additional granularity in $\eta$ and improved uncertainties in the peripheral event range. 
The measurement is also compared to generator-level results from HIJING, EPOS and AMPT to benchmark predictions of MC event generators. 
These measurements are presented across a large centrality range showing a strong understanding of the sPHENIX calorimeter response over at least an order-of-magnitude in dynamic range and making this analysis a necessary step towards accomplishing the sPHENIX jet physics program goals. 